% ============================================================
% Section 136: 液氦温度光激发锗的电流与扩散系数
% ============================================================
\section{半导体物理：低温光激发锗的输运性质}
\label{sec:photoexcited-ge}

\begin{modelbox}[题目：液氦温度下光激发本征锗的电流与扩散]
在液氦温度下，用光激发本征锗产生传导电子，其平均密度为 $1\times10^{12}\,\text{cm}^{-3}$。在此温度下，电子和空穴的迁移率相等，$\mu=0.5\times10^4\,\text{cm}^2/(\text{V}\cdot\text{s})$。

(1) 若将 $100\,\text{V}$ 电压加到边长为 $1\,\text{cm}$ 的立方晶体上，产生的电流约为多少？

(2) 电子和空穴的扩散系数的近似值是多少？（室温下 $k_BT/e=26\,\text{mV}$。）
\end{modelbox}

\begin{tipbox}[考点快速分析]
\begin{itemize}
    \item \textbf{本征光激发}：$n=p$，总电导率为电子和空穴贡献之和
    \item \textbf{电导率}：$\sigma=e(n\mu_n+p\mu_p)=2ne\mu$（迁移率相等时）
    \item \textbf{爱因斯坦关系}：$D/\mu=k_BT/e$，扩散系数与温度成正比
    \item \textbf{液氦温度}：$T=77\,\text{K}$，$k_BT/e\approx6.67\,\text{mV}$
\end{itemize}
\end{tipbox}

\begin{derivationbox}[标准解题过程]
\textbf{(1) 电流的计算}

光激发本征锗中电子和空穴浓度相等：$n=p=1\times10^{12}\,\text{cm}^{-3}$，迁移率相等 $\mu_n=\mu_p=\mu=0.5\times10^4\,\text{cm}^2/(\text{V}\cdot\text{s})$。

电导率：
\[
\sigma=e(n\mu_n+p\mu_p)=2ne\mu.
\]
外加电压 $V=100\,\text{V}$，晶体边长 $d=1\,\text{cm}$，电场强度 $E=V/d=100\,\text{V/cm}$。

电流密度：
\[
j=\sigma E=2ne\mu E.
\]
代入数值：
\[
j=2\times10^{12}\times1.6\times10^{-19}\times0.5\times10^4\times100=0.16\,\text{A/cm}^2.
\]
晶体截面积 $S=1\,\text{cm}^2$，故电流
\begin{equation}
\boxed{I=jS=0.16\,\text{A}=160\,\text{mA}.}
\end{equation}

\textbf{(2) 扩散系数的计算}

爱因斯坦关系：
\[
\frac{D}{\mu}=\frac{k_BT}{e}.
\]
室温（$T_0=300\,\text{K}$）下 $k_BT_0/e=26\,\text{mV}$。液氦温度 $T=77\,\text{K}$ 时：
\[
\frac{k_BT}{e}=26\times10^{-3}\times\frac{77}{300}\approx6.67\times10^{-3}\,\text{V}.
\]
扩散系数：
\[
D=\mu\cdot\frac{k_BT}{e}=0.5\times10^4\times6.67\times10^{-3}\approx33.35\,\text{cm}^2/\text{s}.
\]
\begin{equation}
\boxed{D\approx3.34\times10^{-3}\,\text{m}^2/\text{s}=33.4\,\text{cm}^2/\text{s}.}
\end{equation}
\end{derivationbox}

\begin{formulabox}[答案汇总]
\begin{center}
\begin{tabular}{ll}
\toprule
\textbf{物理量} & \textbf{数值} \\
\midrule
(1) 电流 $I$ & $0.16\,\text{A}=160\,\text{mA}$ \\
(2) 扩散系数 $D$ & $3.34\times10^{-3}\,\text{m}^2/\text{s}=33.4\,\text{cm}^2/\text{s}$ \\
\bottomrule
\end{tabular}
\end{center}
\end{formulabox}

\begin{insightbox}[物理图像]
\begin{itemize}
    \item 光激发本征半导体中电子和空穴成对产生，浓度相等，总电导率为两者之和。
    \item 低温下迁移率较高（$0.5\times10^4\,\text{cm}^2/\text{V}\cdot\text{s}$），即使载流子浓度仅 $10^{12}\,\text{cm}^{-3}$，在 $100\,\text{V}$ 电压下也能产生可观的电流（$160\,\text{mA}$）。
    \item 爱因斯坦关系将扩散系数与迁移率联系起来，是半导体输运理论的基本关系式。低温下扩散系数因 $T$ 降低而减小。
\end{itemize}
\end{insightbox}
